\subsection{Une contribution indirecte mais stratégique}

Les travaux du comité conduisent à poser le cadre de départ suivant :

\begin{quote}
\textbf{L'\ARSeize{} n'a pas vocation à gérer directement les crises, mais à s'y préparer, à accompagner les acteurs concernés dans son champ de compétence et à participer aux retours d'expérience.}
\end{quote}

Ce positionnement conduit à se situer en complément des dispositifs existants, sans chercher à s'y substituer. L'apport possible de l'association doit donc être envisagé au regard des missions exercées par les autorités publiques, en identifiant les domaines où les compétences des auditeurs peuvent utilement renforcer la préparation, la diffusion d'une culture stratégique, l'analyse ou le retour d'expérience.

Cette doctrine de contribution repose sur une limite claire : l'\ARSeize{} ne crée pas une capacité opérationnelle autonome et ne préjuge pas de la disponibilité effective de ses membres pendant une crise. Elle peut en revanche aider à rendre plus lisibles des compétences dispersées, à favoriser des échanges entre milieux professionnels et associatifs, et à produire des analyses utiles dans un cadre institutionnel maîtrisé.

\subsection{Identification de la valeur ajoutée des auditeurs}

% A ETOFFER

Dans un environnement de gestion de crise déjà structuré, il convient de déterminer ce que les membres de l'\ARSeize{} peuvent apporter concrètement et de situer cet apport dans les phases du cycle de crise où il est le plus pertinent.

Cette clarification permet d'éviter un appel trop général, difficilement exploitable par les acteurs publics, tout en distinguant les contributions qui relèvent de l'expérience professionnelle des auditeurs de celles qui peuvent être portées collectivement par le réseau des \AssociationsUnionIHEDN.

Les réflexions convergent vers la réponse suivante : \textbf{cet apport se situe principalement en amont et en aval des crises}.

\begin{proposition}{Positionnement de l'\ARSeize}
L'association ne saurait se substituer aux autorités publiques, ni constituer une chaîne autonome de gestion de crise. Sa contribution se situe d'abord dans la préparation, la diffusion d'une culture stratégique, l'identification des compétences disponibles, puis, après l'événement, dans des analyses.
\end{proposition}

